\documentclass[11pt]{article}
\usepackage[T1]{fontenc}
\usepackage[utf8]{inputenc}
\usepackage{lmodern}
\usepackage[margin=1in]{geometry}
\usepackage{amsmath,amssymb,amsthm,mathtools}
\usepackage{microtype,booktabs,array}
\usepackage[numbers,sort&compress]{natbib}
\usepackage{xurl}
\usepackage[colorlinks=true,allcolors=blue]{hyperref}
\setlength{\emergencystretch}{2em}
\newcommand{\C}{\mathbb C}
\newcommand{\norm}[1]{\lVert #1\rVert}
\newcommand{\W}{W}
\newcommand{\cbase}{1+\sqrt{2}}
\newcommand{\HH}{\mathcal H}
\newcommand{\im}{\operatorname{Im}}
\newcommand{\re}{\operatorname{Re}}
\DeclareMathOperator{\dist}{dist}
\DeclareMathOperator{\conv}{conv}
\DeclareMathOperator{\diag}{diag}
\DeclareMathOperator{\spec}{spec}
\newtheorem{theorem}{Theorem}[section]
\newtheorem{proposition}[theorem]{Proposition}
\newtheorem{lemma}[theorem]{Lemma}
\newtheorem{corollary}[theorem]{Corollary}
\theoremstyle{definition}
\newtheorem{definition}[theorem]{Definition}
\theoremstyle{remark}
\newtheorem{remark}[theorem]{Remark}
\numberwithin{equation}{section}
\title{Finite numerical-range certificates for polynomial matrix computations:\\
a rational prototype with comparative controls}
\author{Robin Gounder\thanks{Correspondence: \href{mailto:robin@gounder.com}{robin@gounder.com}.}\\{\small Vaionex Corporation}}
\date{Research draft, 7 October 2026}
\hypersetup{pdftitle={Finite numerical-range certificates for polynomial matrix computations: a rational prototype with comparative controls},
pdfauthor={Robin Gounder}}
\begin{document}
\maketitle
\begin{abstract}
A numerical-range inequality becomes a usable error certificate only after
its continuous supremum and its numerical inputs have been bounded.
We present a finite certificate for matrix-valued polynomials based on
validated supporting half-planes, a boundary mesh, and coefficient norm
bounds. The certificate is valid with the established complete spectral-set
constant \(1+\sqrt{2}\); a separate hypothesis records the sharp constant
\(2\) asserted in recent preprints and in the OpenAI mathematics release.
We prove an explicit enclosure error estimate, distinguish eigenvalue
validation from eigenpair residuals, and separate input uncertainty from
evaluation roundoff. For Richardson residual powers, the classical numerical-radius power
inequality already supplies an unconditional constant-two bound. The
supremum over a polygon reduces exactly to its vertices, and a rational
certificate can be checked without trigonometric functions or
floating-point logarithms.
Comparative controls show that enclosure geometry can lower a sufficient
budget relative to a single-step induced-norm rate bound, while stronger
exact-power certificates beat the prototype on its featured fixture.
On the coupled \(32\times32\) fixture at tolerance \(10^{-8}\), the
numerical-range certificate gives 31 steps, compared with 11 for direct
powers and 113 for the single-step norm-rate bound. The prototype supports
real-rational scalar polynomials; the theorem permits matrix-valued
polynomials. The experiments compare sufficient certificate budgets and
construction costs. The contribution is a finite certificate specification,
an explicit error decomposition, and reproducible diagnostics built from
established spectral-set and numerical-range methods.
\end{abstract}
\noindent\textbf{Keywords:} numerical range; verified computation;
matrix polynomial; spectral set; Richardson iteration; GMRES.\\
\textbf{MSC 2020:} 15A60, 65F10, 65G20, 47A25.

\section{Introduction and scope}\label{sec:intro}
Polynomial computations occur in iterative linear solvers, matrix-function
approximations, and stability analysis. For nonnormal matrices, evaluating
a polynomial only at the eigenvalues can underestimate its operator norm.
The numerical range retains directional information that can control this
norm. A theorem of Crouzeix and Palencia gives, uniformly in the matrix and
coefficient dimensions,
\begin{equation}\label{eq:CPintro}
 \norm{Q[A]}\le(\cbase)\max_{z\in\W(A)}\norm{Q(z)}
\end{equation}
for matrix-valued polynomials \citep{CrouzeixPalencia2017}.
This complete bound permits coefficient matrices that do not commute.
For an application, however, the right-hand side still contains a
continuous set and a continuous maximization problem. An inaccurate outer
enclosure or an uncertified eigenvalue can defeat the intended guarantee.

The numerical-range functional calculus builds on earlier operator
integral representations \citep{DelyonDelyon1999}, convex-domain
spectral-set estimates \citep{BadeaCrouzeixDelyon2006}, and uniform
numerical-range bounds for polynomial functions
\citep{Crouzeix2007}. Our finite arithmetic interface applies these
established spectral-set estimates to finite, reproducible norm certificates.

The numerical determination of the field of values has a substantial
history. Johnson developed its approximation using directional Hermitian
eigenproblems \citep{Johnson1978}. Loisel and Maxwell developed a
path-following algorithm, with error bounds for specified matrix classes
and numerical cost--accuracy comparisons \citep{LoiselMaxwell2018}.
Field-of-values computation is also available in existing scientific
software, for example the Chebfun exposition and implementation
\citep{Trefethen2010}. These methods supply the supporting-line geometry
used in the present certificate.

Numerical-range estimates already support matrix-function algorithms.
Beckermann and Reichel connect Faber approximation and matrix-function
error estimates \citep{BeckermannReichel2009}. Crouzeix and Greenbaum give
spectral-set estimates with applications to GMRES and rational Krylov
methods \citep{CrouzeixGreenbaum2019}. Our exact-arithmetic GMRES
corollary connects the finite certificate with the standard
residual-minimization property of the method \citep{SaadSchultz1986}.

The polynomial-approximation viewpoint also underlies geometric
analyses of Krylov methods \citep{EiermannErnst2001}.
Beckermann obtained numerical-range GMRES estimates through Faber
polynomials \citep{Beckermann2005}. Eigenvalues alone cannot prescribe
a nonnormal GMRES convergence history: Greenbaum, Pt{\'a}k and
Strako{\v s} constructed any prescribed nonincreasing positive residual
curve with any chosen nonzero spectrum
\citep{GreenbaumPtakStrakos1996}.


The sharp constant is also the subject of literature preceding the
repository studied here. Lorist and Schwenninger assert a general scalar
constant-two proof in a 2026 preprint \citep{LoristSchwenninger2026}.
Crouzeix and Greenbaum subsequently use that result to update
spectral-set estimates \citep{CrouzeixGreenbaum2026Update}. Complete
constant-two results have been reported for weighted shifts
\citep{CrouzeixGreenbaum2025} and for matrices of order at most three
\citep{AhagCzyzVirtanen2026}. The OpenAI release asserts the complete
inequality in arbitrary finite dimensions
\citep{OAI:A-direct-proof-of-the-complete-Crouzeix-inequality-September-26-2026}.
Scalar constant-two applications therefore have a literature basis
preceding the repository release.

This paper addresses a narrower operational question: what finite evidence
suffices to turn a numerical-range bound into a reproducible norm
certificate? Its contributions are an explicit finite-input theorem, an
error decomposition that can be audited separately, and an exact rational
realization for scalar residual polynomials. The geometric ingredients and
the spectral-set applications are established methods. The controls below
test recomputability and quantify conservatism through comparisons with
competing certificates, including norms of matrix powers.

Matrix-function theory and conditioning are treated systematically by
Higham \citep{Higham2008}. Exponential integrators provide a concrete
setting in which matrix functions applied to vectors are required
\citep{HochbruckOstermann2010}; truncated Taylor matrix polynomials
are used in the exponential-action algorithm of Al-Mohy and Higham
\citep{AlMohyHigham2011}. These existing applications motivate future
certificate work. The present prototype implements scalar polynomial and
Richardson certificates; exponential-action algorithms provide a further
application setting.

\paragraph{Evidence status.}
We fix the repository snapshot at commit
\nolinkurl{adc7f1241b42e322a6451854ab7e4b4c146bf78a}, dated
6 October 2026, and record the source audit as of 7 October 2026
\citep{OpenAIRepository2026}. Its scope document describes a sharp complete
constant-two result and links several Lean comparator statements
\citep{OpenAIScope325}. The source assessment uses the supplied scope
document, comparator configuration, and manuscript at this snapshot.
This evidence describes the reported theorem scope and its comparison
statements. For arbitrary polynomials, unconditional conclusions below use
\(\cbase\), and the general constant-two version is stated under an
explicit hypothesis. Fixed Richardson powers instead use the classical
unconditional numerical-radius power inequality, as explained in
Section~\ref{sec:residuals}. This distinction separates the theorem
assumptions from the finite certificate arithmetic.

\section{Notation and the theorem interface}\label{sec:notation}
Let \(A\in\C^{n\times n}\), where \(n\ge1\).
All unsubscripted vector and operator norms are Euclidean norms;
the subscripts \(1\) and \(\infty\) denote the usual induced norms.
Write
\[
 \W(A)=\{x^*Ax:x\in\C^n,\ \norm{x}=1\}.
\]
In finite dimensions the numerical range is nonempty, compact, and convex.
For \(c\in\C\), let \(T=A-cI_n\), and write
\[
 Q(z)=\sum_{k=0}^{d} B_k(z-c)^k,\qquad
 B_k\in\C^{s\times s},\quad s\ge1.
\]
Its operator evaluation is
\begin{equation}\label{eq:evaluation}
 Q[A]=\sum_{k=0}^{d} T^k\otimes B_k.
\end{equation}
The matrix space is the first tensor factor; interchanging the factors is
unitarily equivalent but should not be done silently in software.
Coefficient products are not assumed to commute.

\begin{definition}[Admissible constant]\label{def:constant}
A constant \(C\) is admissible for the computations considered here if
\[
 \norm{Q[A]}\le C\max_{z\in\W(A)}\norm{Q(z)}
\]
holds for every stated matrix and coefficient size and polynomial.
The established universal complete choice is \(C=\cbase\)
\citep{CrouzeixPalencia2017}. The hypothesis \(\HH_2\) is that the same
complete statement holds universally with \(C=2\).
\end{definition}

We use \(\HH_2\) only to label conclusions that require it. For scalar
coefficients, the scalar constant-two statement suffices; the complete
hypothesis is stronger than necessary. Conversely, a scalar theorem alone
does not justify an arbitrary matrix-coefficient certificate.
For a normal \(A\), the spectral theorem supplies \(C=1\) directly:
\[
 \norm{Q[A]}=\max_{\lambda\in\spec(A)}\norm{Q(\lambda)}
 \le \max_{z\in\W(A)}\norm{Q(z)}.
\]
A producer should use such a justified stronger bound when it recognizes
the required structure.

A \emph{validated} quantity in this paper is a proved upper bound for the
specified exact mathematical input. It may be obtained by rational
arithmetic, outward interval arithmetic, or another justified method.
A decimal approximation, an unrecorded numerical tolerance, and a plot are
not certificates. The distinction follows the general philosophy of
verification methods surveyed by Rump \citep{Rump2010}.

\section{Finite outer enclosures}\label{sec:enclosures}
\subsection{Directional support and validated producers}
For \(\theta\in\mathbb R\), set
\[
 H_\theta(T)=\frac{e^{-i\theta}T+e^{i\theta}T^*}{2},\qquad
 h(\theta)=\lambda_{\max}(H_\theta(T)).
\]
For every unit vector \(x\),
\[
 \re(e^{-i\theta}x^*Tx)=x^*H_\theta(T)x.
\]
The Rayleigh--Ritz principle identifies \(h(\theta)\) with the support
function of \(\W(T)\).

Choose directions \(\theta_1,\ldots,\theta_J\) in cyclic order modulo
\(2\pi\), and validated numbers \(\bar h_j\ge h(\theta_j)\). Define
\begin{equation}\label{eq:polygon}
 P=c+\bigcap_{j=1}^{J}
 \{w\in\C:\re(e^{-i\theta_j}w)\le\bar h_j\}.
\end{equation}
This is an outer enclosure. To use its vertices or boundary, its
nonemptiness and boundedness must also be established.

For exact rational geometry, unit normals are unnecessary. Put
\[
 S=\frac{T+T^*}{2},\qquad J_T=\frac{T-T^*}{2i}.
\]
Given a nonzero real normal \(q=(a,b)\), the relevant Hermitian matrix is
\(H_q=aS+bJ_T\), and
\[
 \sup_{z\in\W(T)}\bigl(a\re z+b\im z\bigr)=\lambda_{\max}(H_q).
\]
A validated upper bound \(g_q\) gives the half-plane
\(a\re(z-c)+b\im(z-c)\le g_q\).
With rational \(a,b,g_q,c\), intersections and feasibility checks can be
performed exactly. There is no need to treat rounded sine and cosine values
as exact unit normals.

An elementary valid producer uses the Gershgorin row bound
\begin{equation}\label{eq:gersh}
 g_q=\max_i\left((H_q)_{ii}+\sum_{\ell\ne i} u_{i\ell}\right),
 \qquad u_{i\ell}\ge |(H_q)_{i\ell}|.
\end{equation}
Indeed, for an eigenpair choose a component of maximal modulus. Its
eigenvalue equation gives
\(\lambda\le(H_q)_{ii}+\sum_{\ell\ne i}|(H_q)_{i\ell}|\),
because the eigenvalue is real. This proves the required upper bound.
The producer may be conservative; validity does not imply accuracy.

\begin{remark}[Eigenpair residuals do not certify the upper edge]
For \(H=\diag(0,100)\), the unit vector \(v=(1,0)^T\) has Rayleigh
quotient \(0\) and eigenpair residual \(0\), while
\(\lambda_{\max}(H)=100\). A single small Ritz residual cannot justify an
upper support line. It can assist in a lower bound, or be combined with
additional certified spectral information. Other valid upper producers
include a verified positive-semidefinite test for \(gI-H\) and a validated
full decomposition with controlled orthogonality and remainder.
For a strict upper support, one can instead verify positive definiteness
of \(gI-H_q\). Rump develops rigorous Cholesky-based sufficient tests for
Hermitian positive definiteness with rounding errors included
\citep{Rump2006}. Such a validator is a possible alternative support
producer; it is not supplied by the present implementation.

\end{remark}

\subsection{An explicit enclosure error}
A convergence statement requires more than upper supports. In particular,
refining the directions of a Gershgorin construction need not remove its
spectral slack.

\begin{proposition}[Radius and enclosure excess]\label{prop:geometry}
Assume \(\W(T)\subseteq\{z:|z|\le R\}\), where \(R\ge0\), and let
the maximum cyclic gap between the unit-normal directions be
\(\Delta<\pi\). Suppose, for some \(\eta\ge0\),
\[
 h(\theta_j)\le\bar h_j\le h(\theta_j)+\eta
 \quad\text{for every }j.
\]
Then \(P\) in \eqref{eq:polygon} is a nonempty compact convex outer
enclosure, every \(z\in P\) satisfies
\[
 |z-c|\le\rho,\qquad
 \rho=(R+\eta)\sec(\Delta/2),
\]
and its one-sided Hausdorff excess satisfies
\begin{equation}\label{eq:hausdorff}
 \sup_{z\in P}\dist(z,\W(A))
 \le d_P:=\eta+2(R+\rho)\sin(\Delta/4).
\end{equation}
\end{proposition}
\begin{proof}
The support identity implies \(\W(A)\subseteq P\), proving nonemptiness.
For any nonzero \(w=z-c\) in the intersection, select a sampled normal
within angular distance \(\Delta/2\) of \(w/|w|\). Then
\[
 |w|\cos(\Delta/2)\le
 \re(e^{-i\theta_j}w)\le\bar h_j\le R+\eta.
\]
This establishes boundedness and the radius estimate. The intersection
is closed and convex, hence compact.

Work in shifted coordinates and write \(K=\W(T)\). For \(w\in P-c\)
outside \(K\), let \(y\) be its Euclidean projection onto \(K\).
The unit vector \(u=(w-y)/|w-y|\), identified with a complex unit normal,
supports \(K\) at \(y\). Hence
\(\dist(w,K)=u\cdot w-h_K(u)\).
Choose a sampled unit normal \(u_j\) within angle \(\Delta/2\) of \(u\).
The support function is \(R\)-Lipschitz as a function of the unit normal:
\[
 |h_K(u)-h_K(u_j)|\le R|u-u_j|.
\]
The sampled inequality and \(|w|\le\rho\) give
\[
 u\cdot w-h_K(u)
 \le \eta+(R+\rho)|u-u_j|
 \le \eta+2(R+\rho)\sin(\Delta/4).
\]
Points in \(K\) have distance zero, so the same bound holds throughout.
\end{proof}

The estimate gives an explicit upper budget in terms of support slack
and angular resolution.
Any validated disk bound may serve as \(R\); for instance,
\(R\ge\norm{T}\) is sufficient.
A useful bound on \(\eta\) can be obtained from a certified lower support
value \(\underline h_j\): take
\(\eta=\max_j(\bar h_j-\underline h_j)\).
An exact Rayleigh quotient supplies such a lower value.
The difference must include Gershgorin overestimation as well as arithmetic
rounding. Merely setting \(\eta\) equal to an eigensolver tolerance is
unjustified.

If input uncertainty satisfies \(\norm{T-\widehat T}\le\varepsilon_A\),
a support bound for \(\widehat T\) can be safely enlarged by
\(|q|\varepsilon_A\), since
\(\norm{H_q(T)-H_q(\widehat T)}\le|q|\varepsilon_A\).
This certifies an enclosure of the uncertain input itself. The alternative
nominal-input transfer in Section~\ref{sec:errors} should not be combined
with this enlargement in a way that double-counts the same error.

\section{The finite polynomial certificate}\label{sec:certificate}
Let \(P\) be a nonempty compact convex polygon containing \(\W(A)\).
For a polygon with nonempty interior, let \(Z=\{z_1,\ldots,z_L\}\)
be a finite subset of its boundary such that every boundary point is
within distance \(\delta\) of \(Z\).
If \(P\) is a segment or point, require instead that every point of \(P\)
be within distance \(\delta\) of \(Z\).
This convention avoids invoking a two-dimensional maximum principle on a
degenerate domain.

Choose validated quantities
\[
 \rho\ge\max_{z\in P}|z-c|,\quad
 \beta_k\ge\norm{B_k},\quad
 u_\ell\ge\norm{Q(z_\ell)},
\]
and put
\begin{equation}\label{eq:DandU}
 D_Q=\sum_{k=1}^{d} k\beta_k\rho^{k-1},
 \qquad U_Q=\max_{\ell}u_\ell+\delta D_Q.
\end{equation}
The empty derivative sum is zero. Zeroth powers are interpreted as one,
including when \(\rho=0\).

\begin{theorem}[Finite-input certificate]\label{thm:finite}
Under the stated validated-input conditions,
\[
 \max_{z\in P}\norm{Q(z)}\le U_Q.
\]
Consequently,
\begin{equation}\label{eq:certificate}
 \norm{Q[A]}\le C U_Q
\end{equation}
for every admissible constant \(C\).
The certificate is unconditional for \(C=\cbase\), and its \(C=2\)
version follows under \(\HH_2\).
A rational upper bound \(C_{\rm up}\ge C\) may replace \(C\)
in a numerical certificate.
\end{theorem}
\begin{proof}
On the disk \(|z-c|\le\rho\),
\(\norm{Q'(z)}\le D_Q\).
Integrating along a straight line segment in that convex disk shows that
\(\norm{Q(z)-Q(w)}\le D_Q|z-w|\).
The mesh therefore bounds each covered point by \(U_Q\).

When \(P\) has interior, for any fixed unit vectors \(a,b\in\C^s\),
the scalar polynomial \(a^*Q(z)b\) satisfies the maximum-modulus
principle on \(P\).
For any \(z\in P\),
\[
 |a^*Q(z)b|\le\max_{w\in\partial P}\norm{Q(w)}.
\]
Taking the supremum over \(a,b\) bounds \(\norm{Q(z)}\) by the boundary
supremum, which the mesh estimate bounds by \(U_Q\).
For a degenerate \(P\), the mesh already covers the entire domain.
Finally, \(\W(A)\subseteq P\) and the admissible spectral-set inequality
give \eqref{eq:certificate}.
\end{proof}

\begin{proposition}[Auditable overestimation budget]\label{prop:excess}
Suppose Proposition~\ref{prop:geometry} applies and the sample bounds
also satisfy
\(u_\ell\le\norm{Q(z_\ell)}+\eta_{\rm sample}\), with a known
\(\eta_{\rm sample}\ge0\).
If \(M_Q=\max_{z\in\W(A)}\norm{Q(z)}\), then
\begin{equation}\label{eq:budget}
 U_Q\le M_Q+D_Q(d_P+\delta)+\eta_{\rm sample}.
\end{equation}
\end{proposition}
\begin{proof}
For each sample point choose a closest point in the compact set
\(\W(A)\).
Both points lie in the radius-\(\rho\) disk.
The Lipschitz estimate gives
\(\norm{Q(z_\ell)}\le M_Q+D_Qd_P\).
Adding the sample slack and mesh term proves the claim.
\end{proof}

Theorem~\ref{thm:finite} needs only the upper bounds.
The two-sided accuracy assumptions in Proposition~\ref{prop:excess}
quantify conservatism; they are not needed to establish validity.
Coefficient norm bounds may use a validated Frobenius norm or
\(\sqrt{\norm{B_k}_1\norm{B_k}_\infty}\).
These inexpensive choices can introduce substantial slack for
matrix-valued coefficients.

\paragraph{Finite construction.}
A checker can recompute the supporting inequalities, enumerate the polygon
vertices, bound the enclosing radius, and subdivide each edge.
If an edge of length \(\ell\) is split into \(m\) equal subsegments
and both endpoints are sampled, every point of that edge is within
\(\ell/(2m)\) of a sample.
A validated upper bound on \(\ell\) gives a valid mesh radius.
The point evaluations, coefficient bounds, and all final arithmetic must
be validated.
A producer that checks only the sampled function values has not bounded
the values between samples.

\section{Input uncertainty and evaluation roundoff}\label{sec:errors}
A geometric certificate bounds an exact operator.
If the input is uncertain or an implementation evaluates a nearby
polynomial, those errors require their own evidence.

\begin{proposition}[Nominal-input transfer]\label{prop:perturbation}
Let
\(Q(z)=\sum_{k=0}^d B_k(z-c)^k\) and
\(\widehat Q(z)=\sum_{k=0}^d\widehat B_k(z-c)^k\).
Suppose
\[
 \norm{A-\widehat A}\le\varepsilon_A,\qquad
 \norm{B_k-\widehat B_k}\le\varepsilon_k,\qquad
 \norm{B_k}\le\beta_k,
\]
and
\(\max\{\norm{A-cI},\norm{\widehat A-cI}\}\le M\).
Then
\begin{equation}\label{eq:perturbation}
 \norm{Q[A]-\widehat Q[\widehat A]}
 \le E_{\rm input}
 :=\varepsilon_A\sum_{k=1}^d k\beta_k M^{k-1}
   +\sum_{k=0}^d\varepsilon_k M^k.
\end{equation}
If a finite certificate bounds
\(\norm{\widehat Q[\widehat A]}\le C\widehat U\), it follows that
\(\norm{Q[A]}\le C\widehat U+E_{\rm input}\).
\end{proposition}
\begin{proof}
Write \(T=A-cI\) and \(\widehat T=\widehat A-cI\).
The telescoping identity
\[
 T^k-\widehat T^k
 =\sum_{j=0}^{k-1}T^{k-1-j}(T-\widehat T)\widehat T^j
\]
holds without commutativity and gives
\(\norm{T^k-\widehat T^k}\le kM^{k-1}\varepsilon_A\).
Split each term of the difference as
\[
 T^k\otimes B_k-\widehat T^k\otimes\widehat B_k
 =(T^k-\widehat T^k)\otimes B_k
   +\widehat T^k\otimes(B_k-\widehat B_k).
\]
Use the product rule for the spectral norm of a Kronecker product and
the triangle inequality.
\end{proof}

If an implementation produces \(Y\) with a validated bound
\(\norm{Y-\widehat Q[\widehat A]}\le E_{\rm alg}\), then
\[
 \norm{Y}\le C\widehat U+E_{\rm alg},\qquad
 \norm{Y-Q[A]}\le E_{\rm alg}+E_{\rm input}.
\]
These statements distinguish the computed output norm from its distance to
the intended operator.
The value \(E_{\rm alg}\) requires validation for the evaluation algorithm.
The implementation below supplies exact-rational geometry; floating-point
Horner and Krylov evaluation require their corresponding roundoff analyses.
This separation is consistent with the forward/backward error and
conditioning distinctions in Higham \citep[Chapters~1--3 and~5]{Higham2002}.
An exact polynomial norm certificate therefore does not by itself
provide a floating-point evaluation-error bound.


A similar separation applies to matrix-function approximation.
If \(f\) is holomorphic on a neighborhood of \(P\) and
\(\max_{z\in P}|f(z)-p(z)|\le\varepsilon_{\rm approx}\), the established
functional-calculus bound gives
\[
 \norm{f(A)-p(A)}\le(\cbase)\varepsilon_{\rm approx}.
\]
For a computed approximation \(Y\), a validated evaluation error
\(\norm{Y-p(A)}\le E_{\rm alg}\) gives total error
\((\cbase)\varepsilon_{\rm approx}+E_{\rm alg}\).
This extends the theorem interface to matrix-function approximation.
The scalar reference implementation supplies polynomial mesh bounds;
approximation and evaluation-error bounds remain separate validated inputs.

\section{Residual polynomials and practical decisions}\label{sec:residuals}
\subsection{A vertex certificate for Richardson iteration}
For \(Ax=b\), fixed-step Richardson iteration with parameter
\(\alpha\in\C\) has exact residual
\[
 r_m=(I-\alpha A)^m r_0.
\]
Let the enclosing polygon \(P\) have vertex set \(V\), including endpoints
when \(P\) is a segment and its single point when it is a point.

\begin{lemma}[Vertex reduction]\label{lem:vertices}
Let
\(q_P=\max_{v\in V}|1-\alpha v|\).
For every integer \(m\ge1\),
\[
 \max_{z\in P}|(1-\alpha z)^m|=q_P^m.
\]
\end{lemma}
\begin{proof}
Every point of \(P=\conv(V)\) is a convex combination
\(z=\sum_v t_vv\), with \(t_v\ge0\) and \(\sum_vt_v=1\).
The triangle inequality gives
\[
 |1-\alpha z|
 =\left|\sum_vt_v(1-\alpha v)\right|
 \le\sum_vt_v|1-\alpha v|\le q_P.
\]
The bound is attained at a maximizing vertex.
Raising to the positive integer power \(m\) proves the result.
\end{proof}

Define the numerical radius by \(w(B)=\max_{z\in\W(B)}|z|\),
using complex unit vectors even when \(B\) is real.
The classical numerical-radius power inequality gives
\(w(B^m)\le w(B)^m\); together with \(\norm{B}\le2w(B)\), this yields
\(\norm{B^m}\le2w(B)^m\) \citep[Section~3.2]{Tadmor2025};
see also \citet[Eq.~(1.1)]{LewisOverton2020}.
These established facts give a sharper unconditional bound for this
particular polynomial family than the general Crouzeix--Palencia constant.

\begin{proposition}[Classical power certificate]\label{prop:power}
For \(E=I-\alpha A\) and every integer \(m\ge1\),
\[
 \norm{E^m}\le2w(E)^m\le2q_P^m.
\]
No constant-two Crouzeix hypothesis is needed.
\end{proposition}
\begin{proof}
The affine numerical-range identity gives
\(\W(E)=1-\alpha\W(A)\), so \(w(E)\le q_P\).
Apply the power inequality and \(\norm{E^m}\le2w(E^m)\).
\end{proof}

\begin{corollary}[Residual budget]\label{cor:richardson}
For every \(m\ge1\), the unconditional choice \(C=2\) gives the
following bound. The general spectral-set choice \(C=\cbase\) remains
a valid, weaker comparator:
\[
 \norm{r_m}\le Cq_P^m\norm{r_0}.
\]
If \(0<q_P<1\) and \(0<\tau<1\), a sufficient exact-arithmetic
iteration budget for
\(\norm{r_m}/\norm{r_0}\le\tau\), with \(r_0\ne0\), is
\begin{equation}\label{eq:iterations}
 m_C=\left\lceil
 \frac{\log(\tau/C)}{\log q_P}
 \right\rceil.
\end{equation}
If \(q_P=0\), one step suffices.
If \(q_P\ge1\), this certificate does not guarantee eventual contraction.
\end{corollary}
The constant two here follows from Proposition~\ref{prop:power}, even if
the general hypothesis \(\HH_2\) is unresolved. This distinction is
essential: the Richardson certificate follows from the classical power
inequality. Geometric contractivity is a sufficient certificate condition;
the iteration can converge even when that condition fails.

For real rational \(\alpha\) and rational polygon vertices,
\[
 q_P^2=\max_{v=x+iy\in V}
 \bigl((1-\alpha x)^2+\alpha^2y^2\bigr)
\]
is rational.
A checker can avoid logarithms and square roots in the stopping decision by
testing
\begin{equation}\label{eq:rationalstop}
 C_{\rm up}^2(q_P^2)^m\le\tau^2.
\end{equation}
Testing successive positive integers until this inequality holds gives the
smallest budget certified by the chosen rational constant upper bound.
A configured search cap is a software resource limit, not evidence that
no certificate exists.

\subsection{GMRES and inexact recurrence}
In exact arithmetic, the \(m\)-step GMRES residual minimizes
\(\norm{p(A)r_0}\) over polynomials of degree at most \(m\) with \(p(0)=1\)
\citep{SaadSchultz1986}.
The Richardson polynomial is an admissible comparator. With \(C=2\)
unconditionally by Proposition~\ref{prop:power},
\[
 \norm{r_m^{\rm GMRES}}\le Cq_P^m\norm{r_0}.
\]
This bounds the mathematical minimal residual in exact arithmetic.
Polynomial selection and finite-precision execution require the
corresponding GMRES algorithm analysis.

For a perturbed residual recurrence,
\[
 \widetilde r_{j+1}=(I-\alpha A)\widetilde r_j+\xi_j,\qquad
 \norm{\xi_j}\le\varepsilon_j,
\]
take \(C=2\) unconditionally, define \(\kappa_0=1\) and
\(\kappa_k=Cq_P^k\) for \(k\ge1\).
Unrolling the recurrence proves
\begin{equation}\label{eq:roundoffrecurrence}
 \norm{\widetilde r_m}
 \le Cq_P^m\norm{\widetilde r_0}
    +\sum_{j=0}^{m-1}\kappa_{m-1-j}\varepsilon_j,
 \qquad m\ge1.
\end{equation}
No independence or cancellation of the perturbations is assumed.
If \(q_P<1\) and every \(\varepsilon_j\le\varepsilon\), this gives
\[
 \norm{\widetilde r_m}
 \le Cq_P^m\norm{\widetilde r_0}
 +\varepsilon\left(
 1+Cq_P\frac{1-q_P^{m-1}}{1-q_P}\right).
\]
The quantity \(\varepsilon(1+Cq_P/(1-q_P))\) is a conservative
asymptotic upper bound for the perturbation contribution. It is not a
lower bound on attainable residuals: the permitted perturbations may
all be zero.
Deriving the actual \(\varepsilon_j\), and relating recursive residuals
to explicitly recomputed residuals, remain necessary tasks for a
floating-point solver.
Inexact Krylov methods have their own analysis of controlled
matrix-vector inaccuracies \citep{SimonciniSzyld2003}.
That theory provides algorithm-specific perturbation budgets that can
be used with the present recurrence estimate.


\subsection{When the constant can matter}
For the same validated \(U_Q\), replacing \(\cbase\) by \(2\) reduces
the multiplicative bound by
\[
 1-\frac{2}{1+\sqrt2}=3-2\sqrt2\approx0.1716.
\]
For Richardson powers this replacement is already justified
unconditionally by Proposition~\ref{prop:power}; for a general
polynomial, a constant-two spectral-set conclusion still has its
separate theorem obligation. It does not change the exponential rate
\(q_P\).
For \(0<q_P<1\), the budgets in \eqref{eq:iterations} obey
\[
 0\le m_{\cbase}-m_2
 \le\left\lceil
 \frac{\log((1+\sqrt2)/2)}{-\log q_P}\right\rceil.
\]
The difference may be zero.
These budgets refer to exact constants and exact \(q_P\);
outward numerical bounds can add slack.

The enclosure budget is often more significant.
If \(M_Q>0\) and \(U_Q=M_Q(1+e)\), then the constant-two bound beats
the ideal established bound \((\cbase)M_Q\) exactly when
\[
 e<\frac{\sqrt2-1}{2}\approx0.2071.
\]
This comparison is a diagnostic against an ideal reference, not an
observable acceptance test unless \(M_Q\) is also known.
It explains why a large enclosure excess can erase the advantage of a
sharper constant.
An additive algorithm error can further dominate both bounds.

The alternative induced-norm certificate should always be reported.
For \(E=I-\alpha A\),
\[
 \norm{E^m}\le
 \min\left\{\norm{E}^{m},
 \bigl(\norm{E}_1\norm{E}_\infty\bigr)^{m/2}\right\}.
\]
A validated upper bound on either rate gives a competing budget.
Whenever both it and \(Cq_P^m\) are available, their minimum is a valid
and at least as informative certificate.
The single-step rate is only a basic baseline. It can lose substantial
information that remains in \(E^m\), so beating this baseline is
insufficient evidence of superiority.
Scalable norm estimation also requires care in a validated comparison:
the block one-norm estimator of Higham and Tisseur returns a lower bound,
not an outward upper certificate
\citep[Section~5]{HighamTisseur2000}.
A diagnostic estimate must acquire separate upper-bound validation
before it can replace \(c_j\) in the certificate below.


\begin{proposition}[Standard matrix-power comparison]\label{prop:blocknorm}
Let \(k\ge1\), let \(c_0=1\), and suppose that
\(c_j\ge\norm{E^j}^2\) for \(1\le j\le k\).
For integers \(m=qk+r\), with \(q\ge0\) and \(0\le r<k\),
\[
 \norm{E^m}^2\le c_k^q c_r.
\]
In particular, an exactly computed real rational matrix power gives the
validated choice
\(c_j=\norm{E^j}_1\norm{E^j}_\infty\).
\end{proposition}
\begin{proof}
The identity \(E^{qk+r}=(E^k)^q E^r\) and submultiplicativity give
\(\norm{E^{qk+r}}\le\norm{E^k}^{q}\norm{E^r}\).
Square this inequality and substitute the stated upper bounds.
The induced-norm product bounds the squared Euclidean operator norm.
\end{proof}

This standard comparison supplies the matrix-power baseline.
For an exact rational power, the target test \(c_m\le\tau^2\)
requires no spectral-set constant or square root.
Scanning successive powers can yield a stronger sufficient budget than
both a single-step rate and a numerical-range enclosure.
Direct powers and block-power reuse have their own arithmetic costs,
including fraction growth. A performance comparison must account for
those costs rather than rank methods by iteration budgets alone.
A practical comparison therefore reports bound strength, construction
costs, and reuse of computed powers.

\section{Reference realization and analytic controls}\label{sec:controls}
\subsection{A rational producer and a recomputation contract}
For a real rational matrix \(A\), write
\(S=(A+A^T)/2\) and \(K=(A-A^T)/2\).
The support matrix for \(q=(a,b)\) is
\[
 H_q=aS-ibK.
\]
Its diagonal entries are \(aA_{ii}\), and an off-diagonal modulus is
\(\sqrt{a^2S_{i\ell}^2+b^2K_{i\ell}^2}\).
A rational upper square root gives a fully rational Gershgorin row witness.
At precision \(b_{\rm sq}\), an integer square-root routine can find the
smallest integer \(t\ge0\) for which
\[
 t^2 q_{\rm den}\ge p\,2^{2b_{\rm sq}},
\]
where \(p/q_{\rm den}\ge0\) is the radicand.
Then \(t/2^{b_{\rm sq}}\) is an outward upper bound, with excess less
than \(2^{-b_{\rm sq}}\); exact dyadic roots have zero excess.
This arithmetic statement does not bound the separate Gershgorin slack.

One direction set consists of the integer points on the boundary of
\([-k,k]^2\), \(k\ge1\), with duplicate corners removed.
It contains \(8k\) directions and the four axes.
Axis constraints prove boundedness.
All pairwise nonparallel line intersections are formed rationally and tested
against every inequality.
Their feasible convex hull produces the bounded polygon, including segment
and point cases.
This elementary enumeration makes each feasible intersection auditable
and suits the dimensions and direction counts of the prototype.

A certificate should record the exact input entries and \(\alpha\), the
normal set, all row witnesses, all support bounds, feasible vertices,
\(q_P^2\), the theorem mode, the rational upper constant, and the stopping
inequality.
A checker must recompute these quantities from the input.
It should reject inconsistencies between the declared input and its
numerical evidence, including an inward support, a missing extremal
vertex, an understated squared radius, or a false stopping inequality.
The caller must compare the declared matrix with the intended input
before relying on the result.
Rational data should be serialized as exact numerator--denominator pairs or
unambiguous rational strings, not rounded JSON numbers.

The checker recomputes the arithmetic evidence for the stated exact input
and theorem mode.
The implementation supports scalar polynomials with real rational
coefficients and Richardson polynomials with positive real rational step
parameters, for real rational input matrices. Application use requires
validation of the model that supplied \(A\) and of the solver evaluation.
The general matrix-valued mesh theorem specifies an extension to complex
inputs and matrix coefficients, with coefficient and point-norm validators
as the required additional components.

\subsection{Normal and sharp controls}
For \(A=\diag(1,2)\) and \(\alpha=2/3\), the numerical range is
\([1,2]\) and \(q_P=1/3\) for the exact segment enclosure.
The true residual norm is \(3^{-m}\).
The spectral theorem gives the stronger constant \(1\).
A universal-\(C\) result provides a consistency check; the spectral
bound is stronger for this matrix.

For
\[
 J=\begin{pmatrix}0&2\\0&0\end{pmatrix},
\]
the numerical range is the closed unit disk and \(\norm{J}=2\).
The scalar polynomial \(p(z)=z\) has ratio \(2\).
This is the standard sharpness control; the low-dimensional and
weighted-shift setting was studied before this release
\citep{CrouzeixGreenbaum2025}.
On \(I+J\), Richardson with \(\alpha=1\) has \(q=1\) on the true
numerical range and no strict contraction certificate of the stated form,
even though the residual is zero after two exact steps.
Failure to certify contraction must not be reported as failure of the
iteration.

Let \(N_n\) be the nilpotent upper shift, with ones on the first
superdiagonal.
For \(A=I_n+tN_n\), \(t>0\), and \(\alpha=1\),
\[
 (I-A)^m=(-t)^mN_n^m,\qquad
 \norm{(I-A)^m}=
 \begin{cases}
 t^m,&0\le m<n,\\
 0,&m\ge n.
 \end{cases}
\]
For \(n\ge2\), the induced \(1\)- and infinity-norm rate is \(t\);
for \(n=1\) it is zero. On the displayed \(n=32,t=1/5\) control,
the Gershgorin certificate needs 13 steps versus 12 for the one-step
norm bound. This is not a uniform ordering over the family: for
\(n=2,t=99/100\), eight directions with 40-bit roots certify 31
steps versus 1,833 from the one-step norm rate, while direct powers
certify two steps. Certificate validity and a selected-baseline
improvement are assessed together with the competing exact-power budget.

\subsection{A coupled nonnormal family}\label{sec:coupled}
For even \(n\ge4\), define the real upper-bidiagonal matrix by
\[
 A_{ii}=
 \begin{cases}1,&i\text{ odd},\\11/10,&i\text{ even},\end{cases}
 \qquad
 A_{i,i+1}=
 \begin{cases}-4/5,&i\text{ odd},\\-1/100,&i\text{ even},\end{cases}
\]
and take \(\alpha=1\).
The nonzero weak links join the two-dimensional blocks into a coupled
nonnormal matrix. The displayed rational entries define the analytic fixture.

For \(E=I-A\), direct row and column sums give
\[
 \norm{E}_\infty=4/5,\qquad
 \norm{E}_1=9/10,\qquad
 q_{\rm induced}^2=18/25.
\]
For the real-axis supports of \(A\), the largest off-diagonal Gershgorin
row radius is \(r=(4/5+1/100)/2=81/200\).
The four axis supports alone give the valid rectangle
\[
 P_{\rm box}=
 [\,1-r,\ 11/10+r\,]+i[-r,r].
\]
Its Richardson vertex value is exactly
\begin{equation}\label{eq:coupled}
 q_{\rm box}^2=(1/10+r)^2+r^2
 =\frac{8381}{20000}<\frac{18}{25}<1.
\end{equation}
These statements use rational arithmetic and the proved row bound.
They hold independently of \(\HH_2\).
For a relative target \(\tau=10^{-8}\),
\eqref{eq:rationalstop} with \(C_{\rm up}=17/7\) gives a sufficient box
budget of \(45\) iterations, while the single-step induced-norm rate
budget is \(113\).
Here \(17/7\ge1+\sqrt2\), since \((10/7)^2>2\).
The inequalities were checked exactly; the corresponding \(44\)- and
\(112\)-step tests fail.
These are conservative sufficient budgets for the exact residual bound.

For \(n=32\), \(k=1\), and \(b_{\rm sq}=40\), the eight-direction
prototype gives
\begin{equation}\label{eq:exactq}
 q_P^2=
 \frac{1069775589843180464514493}
      {3777893186295716170956800}
 \approx0.283167240864243.
\end{equation}
Thus \(q_P\approx0.532134607843018\).
The polygon has eight vertices; one maximizing vertex is
\[
 v=\frac{4136912499509}{2748779069440}
    +i\frac{92225107213}{549755813888}.
\]
The other is its complex conjugate.
Using
\(C_{\rm up}=663613970941/274877906944\ge1+\sqrt2\),
the exact test \eqref{eq:rationalstop} gives \(31\) steps.
The same test with \(C_{\rm up}=2\) from the unconditional
numerical-radius power inequality also gives \(31\) steps.
Both \(30\)-step tests fail.
A separate rational recomputation of all supporting row bounds and feasible
line intersections reproduced \eqref{eq:exactq}.
The reduction relative to the single-step \(113\)-iteration rate budget
comes from enclosure geometry. Direct powers provide a stronger competing
certificate, as quantified below.

\paragraph{Stronger exact-power certificate.}
For this same \(n=32\) input, exact rational powers of \(E=I-A\)
give
\begin{equation}\label{eq:power11}
 \norm{E^{11}}_1\norm{E^{11}}_\infty
 =
 \frac{505768790783}{24414062500000000000000000000}
 <10^{-16}.
\end{equation}
Therefore \(11\) iterations suffice for the relative target
\(10^{-8}\) under this standard norm-product certificate.
At \(m=10\), the corresponding product is
\[
 \frac{217552976919}{244140625000000000000000000}
 >10^{-16}.
\]
An exact scan of powers \(1,\ldots,11\) first passes at \(11\).
These statements were checked by a separately implemented dense rational
matrix multiplication. The package reference baseline also uses generic
exact matrix multiplication.
Failure of the \(10\)-step upper-bound test does not prove that an actual
residual cannot reach the target earlier.
The \(11\)-step power certificate beats both numerical-range budgets of
\(31\). This fixture demonstrates recomputability and the conservatism of
the single-step baseline, while direct powers supply the strongest
displayed certificate.

\begin{table}[t]
\centering
\small
\begin{tabular}{lrrrrr}
\toprule
Control & \(q_P\) & \shortstack{One-step\\rate}
 & \shortstack{Exact\\power} & \(C_0\) & \(C_2\)\\
\midrule
Normal \(\diag(1,2)\) & \(0.333333\) & \(17\) & \(17\) & \(18\) & \(18\)\\
\(I_{32}+\frac15N_{32}\) & \(0.216478\) & \(12\) & \(12\) & \(13\) & \(13\)\\
Shifted sharp \(I_2+J\) & \(1.082392\) & --- & \(2\) & --- & ---\\
Coupled family, \(n=32\) & \(0.532135\) & \(113\) & \(11\) & \(31\) & \(31\)\\
\bottomrule
\end{tabular}
\caption{Sufficient relative residual budgets for \(\tau=10^{-8}\).
Here \(C_0=1+\sqrt2\) is implemented with a rational upper bound,
and the \(C_2=2\) column uses the unconditional numerical-radius
power inequality for Richardson powers.
The one-step rate uses
\((\norm{E}_1\norm{E}_\infty)^{m/2}\).
The exact-power column uses
\(\norm{E^m}_1\norm{E^m}_\infty\le\tau^2\).
Both norm certificates have constant one.
Decimals are descriptive; stopping decisions use exact rational
inequalities. Dashes mean that the corresponding geometric or norm-rate
certificate does not certify contraction. Exact powers certify the
shifted sharp control in two steps. The table compares sufficient
certificate budgets.}
\label{tab:controls}
\end{table}

The controls in Table~\ref{tab:controls} test distinct properties:
degenerate geometry, a cheap dominating bound, sharpness with unsuccessful
geometric contraction certification, and a weak one-step baseline exposed
by exact powers.
Where numerical-range budgets are available, the strongest displayed
norm certificate matches or improves them.
The cost study below measures construction, checking and direct powers
on specified analytic fixtures and discretization models.
Validated spectral-norm methods and block-power reuse provide additional
comparators for a broader performance study. Certificate selection depends
on both bound quality and construction cost.

\subsection{A scalar enclosure and mesh error study}\label{sec:ellipse}
To exercise the general-polynomial mesh certificate, rather than the
Richardson vertex shortcut, consider
\[
 A=\begin{pmatrix}-1&1\\0&1\end{pmatrix},\qquad A^2=I.
\]
For an integer normal \((a,b)\), the support Hermitian matrix has eigenvalues
\(\pm\tfrac12\sqrt{5a^2+b^2}\). Thus its numerical range is the ellipse
\(x^2/(5/4)+y^2/(1/4)\le1\).
The polynomials \(p_1(z)=z^2\) and \(p_2(z)=1-z^2\) both have
\(M_p=\max_{\W(A)}|p|=5/4\), whereas
\(\norm{p_1(A)}=1\) and \(\norm{p_2(A)}=0\).
For the second scalar maximum, write the ellipse boundary as
\(z=(\sqrt5\cos\theta+i\sin\theta)/2\). Then
\[
 |1-z^2|^2=25/16-(5/2)t+t^2,\qquad t=\cos^2\theta\in[0,1].
\]
The convex quadratic has maximum \(25/16\) at \(t=0\).
The zero matrix polynomial is a deliberate negative control: a scalar
supremum certificate need not detect annihilation at the matrix.

The study compares the shipped Gershgorin support construction with an
experimental analytic support producer that rounds the displayed support
outward. This experimental producer exercises the support-accuracy
interface; the shipped backend uses Gershgorin bounds. We cross resolutions
\(k=1,2,4,8\), square-root precisions \(b=8,24,40\), and
\(s=4,16,64\) segments per edge, for both producers and polynomials:
144 cases in total. If \(\bar h_q\) is the unnormalized produced support,
set
\[
 R_b=\operatorname{sqrt\_upper}(5/4,b),\qquad
 \eta=\max_q\frac{\bar h_q-
 \max\{0,\operatorname{sqrt\_upper}((5a^2+b_q^2)/4,b)-2^{-b}\}}{k},
\]
where \(q=(a,b_q)\) distinguishes a normal coordinate from the bit precision.
All normals have length at least \(k\), and their angular gap satisfies
\(\Delta\le\arctan(1/k)\le1/k\). Consequently the following rational
relaxations of Proposition~\ref{prop:geometry} suffice:
\[
 \rho=\frac{R_b+\eta}{1-1/(8k^2)},\qquad
 d_P=\eta+\frac{R_b+\rho}{2k},\qquad D_p=2\rho.
\]
They use \(\cos x\ge1-x^2/2\) and \(\sin x\le x\), avoiding unchecked
trigonometric evaluations. The mesh radius \(\delta\) is an outward upper
bound for the longest edge divided by \(2s\). Exact rational assertions
check every vertex radius and all 144 inequalities
\[
 0\le U_p-5/4\le 2\rho(d_P+\delta)+2^{-b}.
\]
The CSV preserves the exact rationals for each term; decimal columns and
Table~\ref{tab:ellipse} are descriptive only.

\begin{table}[t]
\centering\small
\begin{tabular}{lrrrr}
\toprule
Support producer & \(k\) & \(U_{z^2}\) & \(U_{1-z^2}\) & Excess budget\\
\midrule
Analytic & 1 & 1.290326 & 1.278939 & 3.090163\\
Analytic & 2 & 1.271082 & 1.268033 & 1.329167\\
Analytic & 4 & 1.261034 & 1.260257 & 0.642659\\
Analytic & 8 & 1.255562 & 1.255367 & 0.318785\\
Gershgorin & 1 & 2.361888 & 1.426620 & 7.223737\\
Gershgorin & 2 & 2.321651 & 1.343508 & 3.939560\\
Gershgorin & 4 & 2.307298 & 1.313327 & 2.710893\\
Gershgorin & 8 & 2.302639 & 1.304187 & 2.144144\\
\bottomrule
\end{tabular}
\caption{Scalar mesh bounds at 40 root bits and 64 segments per edge.
Both true scalar suprema are \(5/4\). The last column bounds
\(U_p-M_p\), with the same conservative budget for both polynomials.
These are scalar supremum bounds before applying a matrix theorem constant.}
\label{tab:ellipse}
\end{table}

For analytic supports, \(\eta\le2^{-b}/k\); finer directions and meshes
reduce the observed excess. The Gershgorin polygon contains
\(\pm3/2\) for every displayed resolution: its support is
\(|a|+\tfrac12\sqrt{a^2+b_q^2}\), which dominates the support of those points.
Hence its \(z^2\) supremum is at least \(9/4\), leaving an irreducible
excess of at least one over the true \(5/4\), even with arbitrarily many
directions. Mesh refinement alone cannot correct support slack.
The producer-based experiment uses shared geometry helpers and checks
the finite error accounting. The accompanying independent arithmetic
oracle recomputes this study's exact fields using separate root, clipping,
and polynomial-evaluation implementations.

\subsection{Measured certificate costs}\label{sec:cost}
We predefine ten fixtures: the five comparison controls, the
\(n=2,t=99/100\) nilpotent counterexample, the coupled family at \(n=8\),
and three centered-difference matrices for
\(-\nu u''+v u'+r u\) on \((0,1)\) with homogeneous Dirichlet boundaries.
The latter use \(\nu=1/100\), \(v=1/10\), \(r=1\),
\(h=1/(n+1)\), \(n=8,16,32\), diagonal \(d=1+2\nu/h^2\),
lower diagonal \(-\nu/h^2-v/(2h)\), and upper diagonal
\(-\nu/h^2+v/(2h)\). We set \(\alpha=1/d\).
These specified mathematical models supply reproducible discretization
fixtures for the certificate-cost comparison.

Every request uses eight directions, 40 root bits, relative tolerance
\(10^{-8}\), and the unconditional Richardson constant two. Three
sequential observations time each complete producer call, checker call,
and dense exact-power search, after one untimed \(1\times1\) producer
warmup. Construction includes the one-step baseline; checking includes
full regeneration and JSON evidence comparison. The direct-power search
checks every smaller positive power up to the first certificate or the
fixed cap 64. A separate norm-only call normalizes the request, constructs
\(I-\alpha A\), computes its one-step norm product and searches its budget,
using shared arithmetic helpers while skipping all support geometry. Serialization time is excluded from construction; compact
UTF-8 JSON sizes measure evidence size, not peak memory.
The recorded environment is Python 3.11.5 on Windows;
\nolinkurl{results/certificate-cost-study.json} stores its full platform,
all observations and source hash. Timings are environment-specific
medians of certificate construction, checking, and reference calls.

\begin{table}[t]
\centering\footnotesize
\setlength{\tabcolsep}{4pt}
\begin{tabular}{lrrrrrrr}
\toprule
Fixture & NR & One & Power & Build & Check & Norm & Power\\
\midrule
Normal & 18 & 17 & 17 & 1.05 & 1.07 & 0.08 & 0.45\\
Nilpotent \(n=32,t=1/5\) & 13 & 12 & 12 & 85.78 & 85.90 & 7.14 & 69.58\\
Coupled \(n=32\) & 31 & 113 & 11 & 84.36 & 85.10 & 7.16 & 72.92\\
Shifted sharp & --- & --- & 2 & 1.17 & 1.16 & 0.05 & 0.08\\
Slow scalar & 1902 & 1833 & cap & 3.98 & 3.96 & 1.47 & 0.86\\
Nilpotent \(n=2,t=99/100\) & 31 & 1833 & 2 & 2.75 & 2.75 & 1.49 & 0.08\\
Coupled \(n=8\) & 31 & 113 & 11 & 6.20 & 6.24 & 0.54 & 4.12\\
Discretization \(n=8\) & 41 & 39 & 29 & 6.28 & 6.29 & 0.53 & 12.73\\
Discretization \(n=16\) & 121 & 116 & cap & 22.50 & 22.73 & 1.96 & 133.98\\
Discretization \(n=32\) & 427 & 411 & cap & 88.32 & 90.38 & 7.75 & 738.15\\
\bottomrule
\end{tabular}
\caption{Three-observation median certificate timings.
NR and One denote numerical-range and one-step iteration budgets;
the first Power column is an iteration budget. Build, Check, Norm and
the last Power column are median milliseconds. ``cap'' is an unsuccessful
64-power run and must not be ranked as a successful competing certificate.
A dash indicates absence of a geometric contraction certificate.}
\label{tab:cost}
\end{table}

For the coupled \(n=32\) fixture, the exact-power certificate is both
stronger and cheaper to construct than the numerical-range certificate
in these observations. All three discretization cases are already
certified by the cheaper one-step norm route with smaller budgets:
39 versus 41, 116 versus 121, and 411 versus 427, respectively.
The norm-only timing separates that route's work from unnecessary support
construction. At sizes 16 and 32 the direct-power search is inconclusive
at its cap; this does not make numerical-range geometry necessary for
certification. Longer or block-power searches could change both cost
and bound quality. The successful norm-based comparators account for the
displayed ordering of certificate budgets and construction costs.
The numerical-range certificate evidence ranges from 1647 to 15263 bytes
on these fixtures. The measured quantities are certificate timings and
evidence sizes; solver runtime, residual histories, peak memory, and
validated spectral-norm comparisons require a broader study.

\section{Limitations, reproducibility, and research use}\label{sec:limits}
The results certify polynomial operators for the stated exact inputs.
Physical model adequacy and finite-precision solver behavior require
their corresponding model and roundoff analyses. Certificate selection
depends on bound strength and construction cost.
On the featured fixture, the standard exact-power certificate is stronger
than either numerical-range mode.
Coefficient growth can make the derivative estimate large.
A basis other than shifted monomials may reduce this effect, but requires a
corresponding proved derivative bound.

The exact rational scalar prototype is restricted.
Sparse validated eigenvalue methods, scalable half-plane intersection,
general matrix-coefficient sampling, adaptive meshes, and a verified
floating-point solver interface are further engineering tasks.
The enclosure excess estimate requires validated support accuracy if it is
used as a quantitative convergence bound.
Without that evidence, the main certificate can remain valid but
conservative.

The analytic controls are completely specified in
Section~\ref{sec:controls}; they can be independently reconstructed without
repository code or a random seed.
The accompanying \texttt{nrcert} version 0.1.0 source bundle includes the
exact requests, generated certificates, a 30-row comparison CSV, and
command logs. A wheel installed into a separate Python 3.11.5 environment
passed all 44 tests, including evidence mutation and recomputation tests,
and all six command-line examples. The commands are reproduced by
\nolinkurl{reproduce/verify_package.py} in the complete source distribution,
which also contains the test suite; the environment and exit statuses
are recorded in \nolinkurl{audit/installed-package-verification.json}.
Producer and checker share an implementation. The checker recomputes the
complete normalized request and compares regenerated evidence, providing
reproducibility and tamper detection relative to that implementation.
Separate arithmetic validation is supplied by the independent ellipse
oracle for its specified finite study. Together, the exact fixtures,
regeneration checks, and recorded comparisons support the prototype's
certificate specification and error accounting.
Operational limits for this version are dimension 64, degree 128,
direction resolution 8 (64 normals), 128 mesh segments per edge,
iteration search cap 100000, and square-root precision 8--256 bits.
The reference power routine permits at most 128 powers; this study uses 64.
Raw rational strings and derived factors have 8192-character limits, with
8192-bit raw and 12000-bit derived numerator/denominator budgets.
Compact decimal exponents are checked before conversion, including
supported Unicode decimal digits. Matrix dimension, degree, mesh size and
iteration search are capped; exact arithmetic can still be expensive
within those caps. Factored witnesses avoid output expansion; power-comparison
costs still depend on exponents and rational operand sizes. Malformed or
over-budget input raises the documented error. The interface supports
research computation under these resource limits.
Exploratory floating-point norms are recorded as diagnostics.

\section*{Generative AI disclosure}
OpenAI Codex agents contributed substantially to source discovery, literature synthesis, and mathematical exposition.

\begingroup\small
\bibliographystyle{unsrtnat}
\bibliography{references}
\endgroup
\end{document}
